% 無須修改本檔內容，除非校方修改了
% 封面、書名頁、中文摘要、英文摘要、誌謝、目錄、表目錄、圖目錄、符號說明
% 等頁之格式


% make the line spacing in effect
\renewcommand{\baselinestretch}{\mybaselinestretch}
\large % it needs a font size changing command to be effective

% 注意！！此處只是預設值，不需更改此處
% 請更改 my_names.tex 內容
\newcommand\cTitle{論文題目}
\newcommand\eTitle{MY THESIS TITLE}
\newcommand\myCname{王鐵雄}
\newcommand\myEname{Aron Wang}
\newcommand\myStudentID{M9315048}
\newcommand\advisorCnameA{南宮明博士}
\newcommand\advisorEnameA{Dr.~Ming Nangong}
\newcommand\advisorCnameB{李斯坦博士}
\newcommand\advisorEnameB{Dr.~Stein Lee}
\newcommand\advisorCnameC{徐　石博士}
\newcommand\advisorEnameC{Dr.~Sean~Hsu}
\newcommand\univCname{國立台灣科技大學}
\newcommand\univdepartmentEname{Department of Computer Science and Information Engineering}
\newcommand\univEname{National Taiwan University of science and technology}
\newcommand\deptCname{光電工程研究所}
\newcommand\fulldeptEname{Graduate School of Electro-Optical Engineering}
\newcommand\deptEname{Electro-optical Engineering}
\newcommand\collEname{College of Engineering}
\newcommand\degreeCname{碩士}
\newcommand\degreeEname{Master of Science}
\newcommand\cYear{九十四}
\newcommand\cMonth{六}
\newcommand\cDay{十}
\newcommand\eYear{2006}
\newcommand\eMonth{June}
\newcommand\ePlace{Chungli, Taoyuan, Taiwan}


 % user's names; to replace those default variable definitions
% 填入你的論文題目、姓名等資料
% 如果題目內有必須以數學模式表示的符號，請用 \mbox{} 包住數學模式，如下範例
% 如果中文名字是單名，與姓氏之間建議以全形空白填入，如下範例
% 英文名字中的稱謂，如 Prof. 以及 Dr.，其句點之後請以不斷行空白~代替一般空白，如下範例
% 如果你的指導教授沒有如預設的三位這麼多，則請把相對應的多餘教授的中文、英文名
%    的定義以空的大括號表示
%    如，\renewcommand\advisorCnameB{}
%          \renewcommand\advisorEnameB{}
%          \renewcommand\advisorCnameC{}
%          \renewcommand\advisorEnameC{}


% 論文題目 (中文)
\renewcommand\cTitle{
可驗證之安全系統的應用
}

% 論文題目 (英文)
\renewcommand\eTitle{
A Verifiable Secure protocol in a Secure System 
}

% 我的姓名 (中文)
\renewcommand\myCname{王大明}

% 我的姓名 (英文)
\renewcommand\myEname{Da-Ming Wang}

% 我的學號
\renewcommand\myStudentID{M19150001}

% 指導教授A的姓名 (中文)
\renewcommand\advisorCnameA{陳明明\ 博士}

% 指導教授A的姓名 (英文)
\renewcommand\advisorEnameA{Ming-Ming Cheng, Ph.D.}

% 指導教授B的姓名 (中文)
\renewcommand\advisorCnameB{}

% 指導教授B的姓名 (英文)
\renewcommand\advisorEnameB{}

% 指導教授C的姓名 (中文)
\renewcommand\advisorCnameC{}

% 指導教授C的姓名 (英文)
\renewcommand\advisorEnameC{}

% 校名 (中文)
\renewcommand\univCname{國立雲林科技大學資訊工程系}

% 科系 (英文)
\renewcommand\univdepartmentEname{Department of Computer Science and Information Engineering}

% 校名 (英文)
\renewcommand\univEname{National Yunlin University of Science \& Technology \\
Master Thesis
}

% 學位名 (中文)
\renewcommand\degreeCname{碩士}

% 口試年份 (中文、民國)
\renewcommand\cYear{115}

% 口試月份 (中文)
\renewcommand\cMonth{6} 

% 口試年份 (阿拉伯數字、西元)
\renewcommand\eYear{2026} 

% 口試月份 (英文)
\renewcommand\eMonth{June}

\newcommand\itsempty{}
%%%%%%%%%%%%%%%%%%%%%%%%%%%%%%%
%       nyust cover 封面
%%%%%%%%%%%%%%%%%%%%%%%%%%%%%%%
%
%!!!!!!建議直接使用雲科提供的word範本，這邊LaTeX無法調整到一模一樣!!!!!!
\begin{titlepage}

% no page number
% next page will be page 1
\newgeometry{top=3cm, bottom=3cm, left=3cm, right=3cm}
% aligned to the center of the page
\begin{center}
% font size (relative to 12 pt):
% \large (14pt) < \Large (18pt) < \LARGE (20pt) < \huge (24pt)< \Huge (24 pt)
%

\centering
\Large{\univCname}\\
%\vspace{0.5cm}
%\hfill
\Large{\degreeCname 論文}\\
\vspace{0.3cm}
\Large{\univdepartmentEname }\\
\Large{\univEname }\\
%\makebox[8cm][s]{\textbf{\Huge{\degreeCname 論文（初稿）}}}\\ %顯示論文種類 (中文)

\vspace{40pt}
%\large % it needs a font size changing command to be effective
\Large{\cTitle}\\  % 中文題目
%
%\vspace{0.5cm}
%
\Large{\eTitle}\\ %英文題目
\vspace{60pt}
% \makebox is a text box with specified width;
% option s: stretch
% use \makebox to make sure
% 「研究生：」 與「指導教授：」occupy the same width
\Large{\myCname}\\
\Large{\myEname}\\
\vspace{36pt}

\Large 指導教授：  \Large{\advisorCnameA} \\
\Large Advisor:  \Large{\advisorEnameA} \\[1em]
% 判斷是否有 B 教授
\ifx \advisorCnameB \itsempty \else
    & \makebox[4cm][l]{\Large{\advisorCnameB}} \\[1em]
\fi
    
% 判斷是否有 C 教授
\ifx \advisorCnameC \itsempty \else
    & \makebox[4cm][l]{\Large{\advisorCnameC}} \\[1em]
\fi

\vspace{27pt}
\Large{中華民國\cYear 年\cMonth 月}\\
\Large{ \eMonth  }\Large{\ } \Large{\eYear }
%
\end{center}
% Resume the line spacing to the desired setting
\renewcommand{\baselinestretch}{\mybaselinestretch}   %恢復原設定
% it needs a font size changing command to be effective
% restore the font size to normal
\normalsize
\end{titlepage}
%%%%%%%%%%%%%%

%% 從摘要到本文之前的部份以小寫羅馬數字印頁碼
% 但是從「書名頁」(但不印頁碼) 就開始計算
%\setcounter{page}{1}
\pagenumbering{roman}
%\pagenumbering{arabic}
%%%%%%%%%%%%%%%%%%%%%%%%%%%%%%%
%       指導教授推薦書 
%%%%%%%%%%%%%%%%%%%%%%%%%%%%%%%
%
% insert the printed standard form when the thesis is ready to bind
% 在口試完成後，再將已簽名的推薦書放入以便裝訂
% create an entry in table of contents for 推薦書
% 目前送出空白頁
%\newpage{\thispagestyle{empty}\addcontentsline{toc}{chapter}{\nameInnerCover}\mbox{}\clearpage}%
%\newpage

% 判斷是否要浮水印？
\ifx\mywatermark\undefined 
  \thispagestyle{empty}  % 無頁碼、無 header (無浮水印)
\else
  \thispagestyle{EmptyWaterMarkPage} % 無頁碼、有浮水印
\fi

%%%%%%%%%%%%%%%%%%%%%%%%%%%%%%%%%%%%%%%%%%%%%%%%%%%%%%%%%%%%%%%
%%no page number
%% create an entry in table of contents for 書名頁
%\addcontentsline{toc}{chapter}{\nameInnerCover}
%
%
%% aligned to the center of the page
%\begin{center}
%% font size (relative to 12 pt):
%% \large (14pt) < \Large (18pt) < \LARGE (20pt) < \huge (24pt)< \Huge (24 pt)
%% Set the line spacing to single for the titles (to compress the lines)
%\renewcommand{\baselinestretch}{1}   %行距 1 倍
%% it needs a font size changing command to be effective
%%中文題目
%\Large{\cTitle}\\ %%%%%
%\vspace{1cm}
%% 英文題目
%\Large{\eTitle}\\ %%%%%
%%\vspace{1cm}
%\vfill
%% \makebox is a text box with specified width;
%% option s: stretch
%% use \makebox to make sure
%% 「研究生：」 與「指導教授：」occupy the same width
%\makebox[3cm][s]{\large{研 究 生：}}
%\makebox[3cm][l]{\large{\myCname}} %%%%%
%\hfill
%\makebox[2cm][s]{\large{Student: }}
%\makebox[5cm][l]{\large{\myEname}}\\ %%%%%
%%
%%\vspace{1cm}
%%
%\makebox[3cm][s]{\large{指導教授：}}
%\makebox[3cm][l]{\large{\advisorCnameA}} %%%%%
%\hfill
%\makebox[2cm][s]{\large{Advisor: }}
%\makebox[5cm][l]{\large{\advisorEnameA}}\\ %%%%%
%%
%% 判斷是否有共同指導的教授 B
%\ifx \advisorCnameB  \itsempty
%\relax % 沒有 B 教授，所以不佔版面，不印任何空白
%\else
%%共同指導的教授B
%\makebox[3cm][s]{}
%\makebox[3cm][l]{\large{\advisorCnameB}} %%%%%
%\hfill
%\makebox[2cm][s]{}
%\makebox[5cm][l]{\large{\advisorEnameB}}\\ %%%%%
%\fi
%%
%% 判斷是否有共同指導的教授 C
%\ifx \advisorCnameC  \itsempty
%\relax % 沒有 C 教授，所以不佔版面，不印任何空白
%\else
%%共同指導的教授C
%\makebox[3cm][s]{}
%\makebox[3cm][l]{\large{\advisorCnameC}} %%%%%
%\hfill
%\makebox[2cm][s]{}
%\makebox[5cm][l]{\large{\advisorEnameC}}\\ %%%%%
%\fi
%%
%% Resume the line spacing to the desired setting
%\renewcommand{\baselinestretch}{\mybaselinestretch}   %恢復原設定
%\large %it needs a font size changing command to be effective
%%
%\vfill
%\makebox[4cm][s]{\large{\univCname}}\\% 校名
%\makebox[6cm][s]{\large{\deptCname}}\\% 系所名
%\makebox[3cm][s]{\large{\degreeCname 論文}}\\% 學位名
%%
%%\vspace{1cm}
%\vfill
%\large{A Thesis}\\%
%\large{Submitted to }%
%%
%\large{\fulldeptEname}\\%系所全名 (英文)
%%
%%
%\ifx \collEname  \itsempty
%\relax % 沒有學院名 (英文)，所以不佔版面，不印任何空白
%\else
%% 有學院名 (英文)
%\large{\collEname}\\% 學院名 (英文)
%\fi
%%
%\large{\univEname}\\%校名 (英文)
%%
%\large{in Partial Fulfillment of the Requirements}\\
%%
%\large{for the Degree of}\\
%%
%\large{\degreeEname}\\%學位名(英文)
%
%\large{in}\\
%%
%\large{\deptEname}\\%系所短名(英文；表明學位領域)
%%
%\large{\eMonth\ \eYear}\\%月、年 (英文)
%%
%\large{\ePlace}% 學校所在地 (英文)
%\vfill
%\large{中華民國}%
%\large{\cYear}% %%%%%
%\large{年}%
%\large{\cMonth}% %%%%%
%\large{月}\\
%\end{center}
%% restore the font size to normal
%\normalsize
%\clearpage


%%%%%%%%%%%%%%%%%%%%%%%%%%%%%%%%%%%%%%%%%%%%%%%%%%%%%%%%%%%%%%%%%%%%%
%%%%%%%%%%%%%%%%%%%%%%%%%%%%%%%
%       論文口試委員審定書 (計頁碼，但不印頁碼) 
%%%%%%%%%%%%%%%%%%%%%%%%%%%%%%%
%
% insert the printed standard form when the thesis is ready to bind
% 在口試完成後，再將已簽名的審定書放入以便裝訂
% create an entry in table of contents for 審定書
% 目前送出空白頁

%\newpage{\thispagestyle{empty}\addcontentsline{toc}{chapter}{\nameCommitteeForm}\mbox{}\clearpage}%


%%%%%%%%%%%%%%%%%%%%%%%%%%%%%%%
%       中文摘要 
%%%%%%%%%%%%%%%%%%%%%%%%%%%%%%%
%
\newpage
%\thispagestyle{plain}  % 無 header，但在浮水印模式下會有浮水印

% create an entry in table of contents for 中文摘要
\phantomsection
\addcontentsline{toc}{chapter}{\nameCabstract}
\fancypagestyle{plain}{\pagestyle{fancy}}
% aligned to the center of the page
\begin{center}

% font size (relative to 12 pt):
% \large (14pt) < \Large (18pt) < \LARGE (20pt) < \huge (24pt)< \Huge (24 pt)
% Set the line spacing to single for the names (to compress the lines)
\renewcommand{\baselinestretch}{1}   %行距 1 倍
% it needs a font size changing command to be effective
%\large{\cTitle}\\  %中文題目
%\vspace{0.83cm}
% \makebox is a text box with specified width;
% option s: stretch
% use \makebox to make sure
% each text field occupies the same width
%\makebox[1.5cm][s]{\large{學生：}}
%\makebox[3cm][l]{\large{\myCname}} %學生中文姓名
%\hfill
%
%\makebox[3cm][s]{\large{指導教授：}}
%\makebox[3cm][l]{\large{\advisorCnameA}} \\ %教授A中文姓名
%
% 判斷是否有共同指導的教授 B
\ifx \advisorCnameB  \itsempty
\relax % 沒有 B 教授，所以不佔版面，不印任何空白
\else
%共同指導的教授B
\makebox[1.5cm][s]{}
\makebox[3cm][l]{} %%%%%
\hfill
\makebox[3cm][s]{}
\makebox[3cm][l]{\large{\advisorCnameB}}\\ %教授B中文姓名
\fi
%
% 判斷是否有共同指導的教授 C
\ifx \advisorCnameC  \itsempty
\relax % 沒有 C 教授，所以不佔版面，不印任何空白
\else
%共同指導的教授C
\makebox[1.5cm][s]{}
\makebox[3cm][l]{} %%%%%
\hfill
\makebox[3cm][s]{}
\makebox[3cm][l]{\large{\advisorCnameC}}\\ %教授C中文姓名
\fi
%
%\vspace{0.42cm} %2013/06/06
\vspace{0.38cm}
%
%\large{\univCname}\large{\deptCname}\\ %校名系所名
%\vspace{0.83cm}
%\vfill

\makebox[1.5cm][c]{\textbf{摘要}}\\
\end{center}
% Resume the line spacing to the desired setting
\renewcommand{\baselinestretch}{\mybaselinestretch}   %恢復原設定
%it needs a font size changing command to be effective
% restore the font size to normal
\normalsize
%%%%%%%%%%%%%
\thispagestyle{fancy}
在現代電子製造業中，印刷電路板 (PCB) 的品質檢測是確保產品良率的關鍵環節。傳統的人工目檢方式不僅效率低落，且容易因視覺疲勞導致誤判；而現有的自動光學檢測 (AOI) 設備多基於傳統影像處理演算法，對於複雜背景下的微小瑕疵識別率不佳，且難以適應多樣化的缺陷型態。為解決上述問題，本研究提出了一種基於 YOLOv8 架構的改進型物件偵測演算法，旨在提升 PCB 瑕疵檢測的精確度與即時性。

本研究首先針對 PCB 數據集樣本不平衡的問題，採用 Mosaic 資料增強技術進行預處理；其次，在 YOLOv8 的骨幹網路中引入 CBAM 注意力機制 (Convolutional Block Attention Module)，以強化模型對於微小瑕疵特徵的提取能力；最後，利用 EIoU Loss 替換原有的邊界框回歸損失函數，加快收斂速度並提升定位準確度。實驗結果顯示，本方法在公開的 PCB 瑕疵數據集上，平均精確度 (mAP@0.5) 達到了 96.5\%，相較於原始 YOLOv8 模型提升了 3.2\%；同時，檢測速度達到每秒 58 幀 (FPS)。本研究證實了改進後的演算法能有效平衡檢測速度與精度，具備應用於工業流水線即時檢測的潛力。

上文由 GEMINI 生成作為舉例

%%%%%%%%%%%%%%%%%%%%%%%%%%%%%%%
%       英文摘要 
%%%%%%%%%%%%%%%%%%%%%%%%%%%%%%%
%
\newpage
\thispagestyle{plain}  % 無 header，但在浮水印模式下會有浮水印

% create an entry in table of contents for 英文摘要
\phantomsection
\addcontentsline{toc}{chapter}{\nameEabstract}

% aligned to the center of the page
\begin{center}
% font size:
% \large (14pt) < \Large (18pt) < \LARGE (20pt) < \huge (24pt)< \Huge (24 pt)
% Set the line spacing to single for the names (to compress the lines)
\renewcommand{\baselinestretch}{1}   %行距 1 倍
%\large % it needs a font size changing command to be effective
%\large{\eTitle}\\  %英文題目
%\vspace{0.83cm}
% \makebox is a text box with specified width;
% option s: stretch
% use \makebox to make sure
% each text field occupies the same width
%\makebox[2cm][s]{\large{Student: }}
%\makebox[5cm][l]{\large{\myEname}} %學生英文姓名
%\hfill
%
%\makebox[2cm][s]{\large{Advisor: }}
%\makebox[5cm][l]{\large{\advisorEnameA}} \\ %教授A英文姓名
%
% 判斷是否有共同指導的教授 B
\ifx \advisorCnameB  \itsempty
\relax % 沒有 B 教授，所以不佔版面，不印任何空白
\else
%共同指導的教授B
\makebox[2cm][s]{}
\makebox[5cm][l]{} %%%%%
\hfill
\makebox[2cm][s]{}
\makebox[5cm][l]{\large{\advisorEnameB}}\\ %教授B英文姓名
\fi
%
% 判斷是否有共同指導的教授 C
\ifx \advisorCnameC  \itsempty
\relax % 沒有 C 教授，所以不佔版面，不印任何空白
\else
%共同指導的教授C
\makebox[2cm][s]{}
\makebox[5cm][l]{} %%%%%
\hfill
\makebox[2cm][s]{}
\makebox[5cm][l]{\large{\advisorEnameC}}\\ %教授C英文姓名
\fi
%
%\vspace{0.42cm}
%\large{Submitted to }\large{\fulldeptEname}\\  %英文系所全名
%
\ifx \collEname  \itsempty
\relax % 如果沒有學院名 (英文)，則不佔版面，不印任何空白
\else
% 有學院名 (英文)
%\large{\collEname}\\% 學院名 (英文)
\fi
%
%\large{\univEname}\\  %英文校名
%\vspace{0.83cm}
%\vfill
%
\textbf{\large{ABSTRACT}}\\
%\vspace{0.5cm}
\end{center}
% Resume the line spacing the desired setting
\renewcommand{\baselinestretch}{\mybaselinestretch}   %恢復原設定
%\large %it needs a font size changing command to be effective
% restore the font size to normal
\normalsize
%%%%%%%%%%%%%
In the modern electronics manufacturing industry, quality inspection of Printed Circuit Boards (PCBs) is a critical step in ensuring product yield. Traditional manual inspection is inefficient and prone to errors due to visual fatigue. Furthermore, existing Automated Optical Inspection (AOI) equipment, mostly based on traditional image processing algorithms, struggles with low recognition rates for tiny defects in complex backgrounds and lacks adaptability to diverse defect types. To address these issues, this thesis proposes an improved object detection algorithm based on the YOLOv8 architecture to enhance the accuracy and real-time performance of PCB defect detection.

First, this study applies Mosaic data augmentation to address the issue of class imbalance in the PCB dataset. Second, the CBAM (Convolutional Block Attention Module) attention mechanism is integrated into the backbone network of YOLOv8 to strengthen the feature extraction capability for minute defects. Finally, the EIoU Loss function replaces the original bounding box regression loss to accelerate convergence and improve localization accuracy. Experimental results demonstrate that the proposed method achieves a mean Average Precision (mAP@0.5) of 96.5\% on a public PCB defect dataset, representing a 3.2\% improvement over the original YOLOv8 model. Meanwhile, the detection speed reaches 58 Frames Per Second (FPS). This research confirms that the improved algorithm effectively balances detection speed and accuracy, demonstrating significant potential for real-time inspection in industrial assembly lines.

上文由 GEMINI 生成作為舉例

%%%%%%%%%%%%%%%%%%%%%%%%%%%%%%%
%       誌謝 
%%%%%%%%%%%%%%%%%%%%%%%%%%%%%%%
%
% Acknowledgment
\newpage
\phantomsection
\addcontentsline{toc}{chapter}{\nameAckn}

\chapter*{\protect\makebox[5cm][c]{\large{\nameAckn}}} %\makebox{} is fragile; need protect

本論文得以順利完成，首先要誠摯感謝我的指導教授 [教授姓名] 博士。在 [碩士/博士] 求學期間，老師不僅在學術研究上給予我嚴謹的指導與啟發，更在待人處事上樹立了榜樣。老師對於 [提及一個具體的研究方向或細節，例如：演算法優化/系統架構] 的獨到見解，使我能在研究瓶頸中找到突破口，在此致上最深的謝意。

同時，感謝口試委員 [委員姓名] 教授、[委員姓名] 教授與 [委員姓名] 教授在口試期間提出的寶貴建議與指正，補足了本論文在 [理論推導/實驗驗證] 上的不足，使本論文的架構更加嚴謹完善。

感謝實驗室的 [學長姐姓名] 在研究初期的引領，以及同窗 [同學姓名] 與學弟妹們在 [實驗室生活/專案合作] 中的互相砥礪與協助。那些一起熬夜跑數據、討論程式碼的日子，將是我求學生涯中珍貴的回憶。

最後，謹以此文獻給我摯愛的家人。感謝父母 [父親姓名]、[母親姓名] 多年來無條件的支持與栽培，讓我能無後顧之憂地專注於學業。你們的鼓勵是我前進的最大動力。

上文由GEMINI生成作為舉例






%%%%%%%%%%%%%%%%%%%%%%%%%%%%%%%
%       目錄 
%%%%%%%%%%%%%%%%%%%%%%%%%%%%%%%
%
% Table of contents
\newpage
\phantomsection % 放船錨與寫入目錄，必須要在 \tableofcontents 之前！
\addcontentsline{toc}{chapter}{\nameToc}

% 1. 先把真正的粗體指令備份為 \realbfseries
\let\realbfseries\bfseries
% 2. 在標題中使用備份的粗體指令
\renewcommand{\contentsname}{\protect\makebox[1cm][s]{\protect\realbfseries\nameToc}}
{ 
  % 1. 開啟一個群組 (大括號)，限制修改範圍只在目錄頁有效
  
  % 2.這行指令的意思是：將「粗體指令 (\bfseries)」重新定義為「一般細體 (\mdseries)」
  % 這樣當 LaTeX 試圖把章節變粗時，它會變成一般字體。
  \let\bfseries\mdseries 
  
  % 3. 產生目錄
  \tableofcontents
  
}


%%%%%%%%%%%%%%%%%%%%%%%%%%%%%%%
%       表目錄 
%%%%%%%%%%%%%%%%%%%%%%%%%%%%%%%
%
% List of Tables
\newcommand{\lotlabel}{表}
\newpage
\phantomsection
\addcontentsline{toc}{chapter}{\nameLot}

\renewcommand{\numberline}[1]{\lotlabel~#1\hspace*{1em}}
\renewcommand{\listtablename}{\protect\makebox[1.5cm][s]{\protect\realbfseries\nameLot}}
{ \let\bfseries\mdseries 
  \listoftables 
}
%\makebox{} is fragile; need protect


%%%%%%%%%%%%%%%%%%%%%%%%%%%%%%%
%       圖目錄 
%%%%%%%%%%%%%%%%%%%%%%%%%%%%%%%
%
% List of Figures
\newcommand{\loflabel}{圖}
\newpage
\phantomsection
\addcontentsline{toc}{chapter}{\nameTof}

\renewcommand{\numberline}[1]{\loflabel~#1\hspace*{1em}}
\renewcommand{\listfigurename}{\protect\makebox[1.5cm][s]{\protect\realbfseries\nameTof}}
{ \let\bfseries\mdseries 
  \listoffigures
}
%\makebox{} is fragile; need protect


%%%%%%%%%%%%%%%%%%%%%%%%%%%%%%%
%       符號說明 
%%%%%%%%%%%%%%%%%%%%%%%%%%%%%%%
%
% Symbol list
% define new environment, based on standard description environment
% adapted from p.60~64, <<The LaTeX Companion>>, 1994, ISBN 0-201-54199-8
\newcommand{\SymEntryLabel}[1]%
{\makebox[3cm][l]{#1}}

\newenvironment{SymEntry}
  {\begin{list}{}%
      {\renewcommand{\makelabel}{\SymEntryLabel}%
       \setlength{\labelwidth}{3cm}%
       \setlength{\leftmargin}{\labelwidth}%
       }%
  }%
  {\end{list}}
%
\newpage
\phantomsection
\addcontentsline{toc}{chapter}{\nameSlist}

\chapter*{\protect\makebox[5cm][c]{\nameSlist}} %\makebox{} is fragile; need protect
\input{frontpages/my_symbols.tex}


\newpage
\setcounter{page}{1}
\pagenumbering{arabic}
%% 論文本體頁碼回復為阿拉伯數字計頁，並從頭起算
%\pagenumbering{arabic}
%%%%%%%%%%%%%%%%%%%%%%%%%%%%%%%%